\begin{decisione}{D0004}{Come riconoscere e contare gli ordini doppi}{Proposta}{23/10/2026}

\contesto
Con la coda lo stesso ordine può essere consegnato due volte (D0002). Va salvato una volta sola, ma i doppi vanno anche contati, perché sono una delle misure. Contarli nella memoria del programma non funziona: si azzera proprio quando il programma si spegne, cioè durante i guasti.

\scelta
\begin{itemize}
  \item Ogni ordine ha un \textbf{identificativo scelto dal cliente} (un UUID), che il database accetta una sola volta (chiave primaria).
  \item Il salvataggio usa ``inserisci, e se esiste già non fare nulla'' (\texttt{INSERT \ldots\ ON CONFLICT DO NOTHING}). Se non ha inserito nulla, l'ordine è un doppio: lo registriamo nella tabella \texttt{duplicati} (identificativo, istante, chi l'ha ricevuto), nella stessa operazione.
  \item Al doppio si risponde come al primo (201 o conferma alla coda): per il cliente l'ordine è salvato, ed è vero.
\end{itemize}

\conseguenze
Nessun ordine salvato due volte, e un conteggio dei doppi che sopravvive ai guasti. Lo stesso meccanismo vale in tutte le versioni: in V1, V2 e V4 i doppi dovrebbero essere zero, e trovarne uno indicherebbe un errore nel banco di prova.

\end{decisione}
